\section{Uvod}

Mjerenje pomaka i deformacija na površini opterećenih konstrukcijskih elemenata
temeljni je zadatak eksperimentalne mehanike. Klasične mjerne metode, poput
tenzometarskih traka, daju informaciju samo u pojedinačnim točkama i zahtijevaju
fizički kontakt s uzorkom. Optičke, bezkontaktne metode -- prije svega
\emph{korelacija digitalne slike} (engl. \textit{Digital Image Correlation}, DIC)
-- omogućuju mjerenje cijelog polja pomaka i deformacija na površini uzorka bez
kontakta, uz visoku prostornu razlučivost.

Dvodimenzijska (2D) korelacija digitalne slike koristi jednu kameru i pretpostavlja
da se gibanje odvija u ravnini paralelnoj s ravninom senzora. Ta pretpostavka
narušava se čim se pojavi izvanravninski pomak: površina koja se približava ili
udaljava od kamere, savija ili rotira, unosi prividnu deformaciju koju 2D metoda
ne može razlučiti od stvarne. \emph{Stereo korelacija digitalne slike}
(stereo-DIC ili 3D-DIC) rješava taj problem korištenjem dviju kalibriranih kamera:
kombinacijom stereovizije i korelacije rekonstruira se \emph{trodimenzijsko} polje
pomaka, uključujući i ravninske i izvanravninske komponente.

\subsection{Cilj rada}

Cilj ovog seminarskog rada jest razviti funkcionalan sustav za mjerenje
trodimenzijskih pomaka i deformacija površine objekta pomoću stereo korelacije
digitalne slike, u skladu s postavljenim zadatkom. Konkretno, potrebno je:
\begin{itemize}
  \item postaviti i kalibrirati stereovizijski sustav (dvije kamere);
  \item provesti korelaciju referentne i deformiranih slika za svaku kameru;
  \item rekonstruirati trodimenzijske pomake točaka na površini uzorka;
  \item odrediti pripadno polje deformacija te vizualizirati deformirano stanje.
\end{itemize}

\subsection{Struktura rada}

U poglavlju \ref{sec:teorija} iznesena je teorijska podloga: model kamere,
projekcijske matrice, princip korelacije te stereo triangulacija. Poglavlje
\ref{sec:metoda} opisuje primijenjenu metodologiju i eksperimentalni postav, a
poglavlje \ref{sec:implementacija} programsku implementaciju u MATLAB-u.
Rezultati mjerenja na stvarnom vlačnom pokusu prikazani su i raspravljeni u
poglavlju \ref{sec:rezultati}, nakon čega slijedi zaključak.
